% One-time imported draft. Editable native source; no Markdown regeneration.
\section{Step 10 Print unsigned
decimal}\label{004-bit-inspector-s10-step-10-print-unsigned-decimal}

Service 36 treats that same low word as an unsigned integer. It does not
need a separate positive 64-bit argument to print the unsigned word
interpretation. The format service supplies the interpretation.

\subsection{Write or
inspect}\label{004-bit-inspector-s10-write-or-inspect}

Append the unsigned line.

\MARSCodeLabel{MIPS assembly}
\begin{minted}{mars-idea}
    la $a0, unsigned_msg
    li $v0, 4
    syscall
    move $a0, $s0
    li $v0, 36
    syscall
    la $a0, newline
    li $v0, 4
    syscall
unsigned_printed:
\end{minted}

\subsection{Expected
observation}\label{004-bit-inspector-s10-expected-observation}

The second line is unsigned: 4294967295 for the same default bits that
printed as signed -1.

\subsection{Check your
understanding}\label{004-bit-inspector-s10-check-your-understanding}

Why can the same \texttt{\$s0} argument serve both decimal reports?

\begin{itemize}
\tightlist
\item
  \textbf{A} Service 36 first changes the sample in memory.
\item
  \textbf{B} An unsigned interpretation always flips the sign bit.
\item
  \textbf{C} The services interpret its low word differently.
\end{itemize}
